% ------------------------------------------------------------------------
% Clasificacion del Riesgo de Inundacion por Parroquia - Esmeraldas
% Proyecto 2P - Universidad de las Fuerzas Armadas ESPE
% Compilar con: pdflatex informe.tex
% ------------------------------------------------------------------------
\documentclass[12pt,a4paper]{article}

% ---- Paquetes ----
\usepackage[utf8]{inputenc}
\usepackage[T1]{fontenc}
\usepackage[spanish]{babel}
\usepackage{geometry}
\geometry{margin=2.5cm, includefoot, footskip=30pt}
\usepackage{graphicx}
\usepackage{booktabs}
\usepackage{array}
\usepackage{hyperref}
\hypersetup{
    colorlinks=true,
    linkcolor=blue,
    urlcolor=blue,
    citecolor=blue
}
\usepackage{setspace}
\onehalfspacing
\usepackage{fancyhdr}
\pagestyle{fancy}
\fancyhf{}
\fancyhead[L]{Riesgo de Inundacion - Esmeraldas}
\fancyhead[R]{\thepage}
\renewcommand{\headrulewidth}{0.4pt}

% ---- Titulo ----
\title{\vspace{-1.5cm}
    \textbf{Clasificacion del Riesgo de Inundacion por Parroquia}\\
    \large Provincia de Esmeraldas, Ecuador
}
\author{
    \textbf{Proyecto 2P -- Aprendizaje Automatico} \\
    Universidad de las Fuerzas Armadas ESPE \\
    \texttt{github.com/tu-usuario/riesgo-inundacion-esmeraldas}
}
\date{\today}

\begin{document}
\maketitle
\thispagestyle{empty}

% ------------------------------------------------------------------------
\section{Introduccion y Contextualizacion}
% ------------------------------------------------------------------------

Ecuador presenta una elevada vulnerabilidad frente a eventos hidrometeorologicos extremos, particularmente inundaciones, que se intensifican durante las temporadas de lluvia en la region Costa. La provincia de Esmeraldas, ubicada en el noroccidente del pais, registra precipitaciones anuales que oscilan entre 1900 y 3000 mm, con zonas de topografia plana y cercania a cuencas hidrograficas que incrementan su susceptibilidad.

La parroquia constituye la unidad territorial minima con informacion demografica, geografica y climatica suficiente para el analisis de riesgo con enfoque local. Este proyecto disena e implementa un modelo de clasificacion supervisada que asigna a cada parroquia una categoria de riesgo de inundacion --- Bajo, Medio o Alto --- integrando variables climaticas, topograficas y socioterritoriales a partir de fuentes de datos abiertas y oficiales.

% ------------------------------------------------------------------------
\section{Datos y su Origen}
% ------------------------------------------------------------------------

Los datos provienen de la Secretaria de Gestion de Riesgos Ecuador (SNGRE) a traves de su REST API publica. Se extrajeron tres capas fundamentales:

\begin{itemize}
    \item \textbf{ZONAS\_SUSCEPTIBLES\_INUNDACIONES:} 320 poligonos de susceptibilidad con atributos de nivel de inundabilidad (ALTA, MEDIA, BAJA) y area en hectareas.
    \item \textbf{COE2:} 30 registros historicos de eventos de inundacion con detalle de afectados, viviendas danadas y ubicacion por parroquia.
    \item \textbf{LIMITE\_PARROQUIAL:} 64 poligonos con los limites oficiales de las parroquias de Esmeraldas.
\end{itemize}

La variable objetivo se construyo combinando un puntaje de susceptibilidad (ponderacion del area con nivel ALTA y MEDIA) con los eventos historicos (numero de inundaciones y total de afectados). Se definieron los umbrales: $\geq$ 35 puntos como Alto, $\geq$ 15 como Medio y menor a 15 como Bajo.

Las variables predictoras --- precipitacion anual, altitud media, pendiente del terreno, distancia a cuerpos de agua, densidad poblacional y area urbanizada --- se generaron con distribuciones diferenciadas segun el nivel de riesgo, manteniendo valores realistas y con correlacion controlada con la variable objetivo.

% ------------------------------------------------------------------------
\section{Metodologia}
% ------------------------------------------------------------------------

El analisis exploratorio incluyo histogramas de distribucion por clase, matriz de correlacion y verificacion de valores nulos. Se observo una separacion clara entre las tres clases en las variables \texttt{DISTANCIA\_AGUA\_KM}, \texttt{PENDIENTE\_PCT} y \texttt{PRECIPITACION\_ANUAL\_MM}, mientras que \texttt{ALTITUD\_MEDIA\_M} y \texttt{DENSIDAD\_POBLACIONAL} mostraron menor poder discriminante.

Se implementaron cuatro modelos base: Regresion Logistica, Arbol de Decision, Random Forest y SVM con kernel RBF. Cada modelo se integro en un Pipeline con RobustScaler para normalizacion robusta ante valores atipicos. Se utilizo validacion cruzada estratificada de 5 pliegues con metrica F1-weighted. Adicionalmente, se implementaron dos ensambles Voting (Hard y Soft) que combinan todos los clasificadores base.

La optimizacion de hiperparametros se realizo mediante GridSearchCV sobre Random Forest, explorando combinaciones de n\_estimators (100, 200, 300), max\_depth (5, 8, 12, None), min\_samples\_split (2, 4, 6) y min\_samples\_leaf (1, 2, 4). El mejor modelo se exporto como Pipeline completo con Joblib.

% ------------------------------------------------------------------------
\section{Resultados}
% ------------------------------------------------------------------------

\subsection{Comparacion de Modelos}

\begin{center}
\begin{tabular}{lccc}
\toprule
\textbf{Modelo} & \textbf{Accuracy} & \textbf{F1-Score} & \textbf{CV F1} \\
\midrule
Random Forest (opt) & 92.71\% & 92.68\% & 93.27\% \\
Voting Soft & 92.71\% & 92.68\% & 93.21\% \\
Logistic Regression & 91.67\% & 91.61\% & 91.41\% \\
SVM (RBF) & 91.67\% & 91.60\% & 91.98\% \\
Decision Tree & 88.54\% & 88.56\% & 84.40\% \\
\bottomrule
\end{tabular}
\end{center}

Random Forest optimizado alcanzo el mejor desempeno con Accuracy de 92.71\% y F1-Score de 92.68\%. La validacion cruzada confirmo estabilidad con CV F1 de 93.27\% y desviacion estandar de solo 1.78 puntos porcentuales.

\subsection{Importancia de Caracteristicas}

El analisis de importancia del Random Forest revelo que las tres variables mas influyentes son: distancia a cuerpos de agua (46.6\%), pendiente del terreno (22.9\%) y precipitacion anual (14.9\%). Estos resultados son coherentes con la literatura sobre riesgo de inundaciones en zonas costeras, donde la proximidad a rios y la topografia plana son factores determinantes.

\subsection{Distribucion de Riesgo}

La aplicacion del modelo sobre las 64 parroquias de Esmeraldas clasifico: 30 parroquias como Bajo riesgo, 23 como Medio y 11 como Alto. Esta distribucion refleja que, si bien la mayoria del territorio presenta condiciones moderadas, existe un grupo significativo de parroquias con riesgo elevado que requieren atencion prioritaria.

% ------------------------------------------------------------------------
\section{Discusion y Limitaciones}
% ------------------------------------------------------------------------

El modelo desarrollado alcanza un rendimiento elevado, pero presenta limitaciones importantes. Las variables predictoras fueron generadas sinteticamente a partir de distribuciones diferenciadas por nivel de riesgo, lo que constituye una idealizacion de la realidad. En un escenario operativo, estas variables deberian obtenerse de fuentes primarias: estaciones meteorologicas del INAMHI, modelos digitales de elevacion del IGM y datos censales del INEC.

El tamano del dataset (320 registros) es limitado para tareas de clasificacion multiclase. La incorporacion de datos temporales --- precipitacion mensual historica, registros de eventos por temporada --- podria mejorar la capacidad predictiva y permitir modelos de riesgo estacional.

Otra limitacion es que el modelo asigna un riesgo uniforme a toda la parroquia, cuando en realidad el riesgo puede variar significativamente dentro del territorio parroquial. La resolucion espacial podria mejorarse trabajando a nivel de recintos o zonas censales.

% ------------------------------------------------------------------------
\section{Conclusiones}
% ------------------------------------------------------------------------

Se diseno e implemento un modelo de clasificacion supervisada para determinar el nivel de riesgo de inundacion por parroquia en la provincia de Esmeraldas, Ecuador. El modelo Random Forest optimizado mediante GridSearchCV alcanzo un Accuracy de 92.71\% y F1-Score de 92.68\%, con capacidad de discriminar correctamente las tres categorias de riesgo.

Las variables mas influyentes --- distancia a cuerpos de agua, pendiente del terreno y precipitacion anual --- son consistentes con los factores que la literatura especializada identifica como determinantes del riesgo de inundacion.

El modelo se integro en una aplicacion web Flask con mapa interactivo Leaflet, desplegada en produccion, que permite visualizar los resultados por parroquia y realizar predicciones personalizadas. Esta herramienta constituye un insumo util para la toma de decisiones en gestion de riesgos a nivel local.

\bigskip
\noindent\textbf{Palabras clave:} clasificacion supervisada, riesgo de inundacion, Random Forest, Esmeraldas, SNGRE, Flask, Leaflet.

\end{document}
